\section{Intro to log geometry (Tingting Fang and Zijian Han)}%
\label{sec:Intro to log geometry (Tingting Fang and Zijian Han)}

\subsection{Motivation}%
\label{sub:Motivation}

In enumerative geometry, we will consider intersection numbers on compactified moduli spaces. This is going to require our objects to degenerate in some way, for example for domain curves to be nodal, maps to land in boundary divisors, or we degenerate the target. In ordinary geometry, we view these as singular, but in log geometry, there is a way to think of these situations as being \textit{log smooth}. In log geometry, we are going to replace functions with the data of functions and how they vanish on the boundary (or namely the morphism \(\alpha \colon M_X \to \mc{O}_X\)).

\begin{exm}
    Consider a family of curves \((xy=t)\) over \(\A^1\), where the central fiber is nodal. In log geometry, if we place the divisorial log structure on \(\A^1\), then this family is log smooth.
\end{exm}

\subsection{Monoids}%
\label{sub:Monoids}

\begin{defn}
    A \textit{monoid} is a commutative semigroup with a unit. 
\end{defn}

For a monoid \(P\), we may associate a group \(P^{\mr{gp}}\) via the standard Grothendieck construction
\[ P^{\mr{gp}} = \ab\{(a,b) \mid (a,b) \sim (c,d) \text{ if there exists } s \in P \text{ such that } a+d+s = b+c+s \}. \]
\begin{defn}
If the natural morphism \(P \hookrightarrow P^{\mr{gp}}\) is injective, we say that \(P\) is \textit{integral}.
\end{defn}

\begin{exm}
    If \(P\) has generators \(a,b,c\) with relation \(a+c=b+c\), then in \(P^{\on{gp}}\) we have \(a = b\), so \(P\) is not integral.
\end{exm}

\begin{defn}
    A monoid \(P\) is 
    \begin{itemize}
        \item \textit{fine} if \(P\) is finitely generated and integral;
        \item \textit{saturated} if for all \(p \in P^{\mr{gp}}\) and all natural numbers \(n\), if \(np \in P\) then \(p \in P\);
        \item \textit{torsion-free} if \(P^{\mr{gp}}\) is;
        \item \textit{sharp} if the identity is the only invertible element.
    \end{itemize}
\end{defn}

\subsection{Log structures}%
\label{sub:Log structures}

Let \(X\) be a scheme.
\begin{defn}
    A \textit{pre-log structure} on \(X\) is a sheaf of monoids \(M_X\) on the small \'etale site of \(X\) and a morphism \(\alpha \colon M_X \to \mc{O}_X\). It is a \textit{log structure} if \(\alpha^{-1}(\mc{O}_X^{\times}) \simeq \mc{O}_X^{\times}\), where the isomorphism is given by \(\alpha\). 
\end{defn}

We now consider the exact sequence
\[ 0 \to \alpha^{-1}(\mc{O}_X^{\times}) \to M_X \to \bar{M}_X \to 0, \]
and the quotient here is called the \textit{characteristic sheaf} or the \textit{ghost sheaf}.

\begin{exm}
    Let \(M_X = \mc{O}_X^{\times}\)  and \(\alpha\) be the identity. Then there is no additional information provided by the log structure.
\end{exm}

\begin{exm}
    Let \(X\) be a regular variety and \(D\) be a closed subscheme. Set \(U = X \setminus D\). Then define
    \[ M_X = \mc{O}_X \cap j_* \mc{O}_U^{\times}, \]
    where \(j \colon U \to X\) is the open immersion. If \(D\) is a divisor, then this is called the \textit{divisorial log structure}.

    For example, if \(X = \A^1\) and \(D = \ab\{0\}\), then the characteristic sheaf vanishes everywhere except for the origin. Here, we compute
    \[ M_X = \ab\{ \varphi \cdot t^n \mid \varphi \in \mc{O}_X^{\times} \}. \]
    The stalk is \(\mc{O}_X^{\times}\) everywhere except at the origin, where it is \(\mc{O}_{\A^1, 0}^{\times} \oplus \N\).
\end{exm}

\begin{exm}
    Now let \(X = \A^2\) and \(D\) be the union of the coordinate axes. Then
    \begin{align*}
        M_{X,D} &= \ab\{f \in k[x,y] \mid f \text{ is invertible outside } D \} \\
        &= \ab\{\varphi \cdot x^a \cdot y^b\}.
    \end{align*}
    Note here that at the general point, we see \(M_p = \mc{O}_{X,p}^{\times}\). At the origin, we see \(\mc{O}_{X,p}^{\times}\oplus \N^2\). At a general point on the \(x\)-axis, we see that \(x^a\) is invertible, so we see the stalk \(\mc{O}_{X,p}^{\times}\oplus \N \cdot b\). Similarly, we can compute the stalk at a general point on the \(y\)-axis.   
\end{exm}

Just as we can associate a sheaf to a presheaf, we will associate a log structure to a pre-log structure. First, a pushout in a category \(\mc{C}\) will be denoted by \(B \sqcup_A C\) and is the colimit of the diagram
\begin{equation*}
\begin{tikzcd}
    A \ar{r}{f} \ar{d}[swap]{g} & B \\
    C.
\end{tikzcd}
\end{equation*}
In the case where \(\mc{C}\) is commutative monoids, then
\[ B \sqcup_A C = (B \oplus C)/\sim, \qquad (f(a) + b, c) \sim (b, g(a)+c). \]

\begin{exm}
    Consider the morphisms \(2 \colon \N \to \N\) and \(3 \colon \N \to \N\). Then the pushout is the monoid
    \[ P = \frac{\N \oplus \N}{(3, 0) \sim (0,2)} \simeq \N \setminus \{1\}. \]
\end{exm}

If we replace \(\mc{C}\) by the category of sheaves of commutative monoids, we can form the pushout on sections and then sheafify.

\begin{defn}
    The \textit{associated log structure} of a pre-log structure \(\alpha \colon M_X \to \mc{O}_X\) is the pushout 
    \[ M_X \sqcup_{\alpha^{-1}(\mc{O}_X^{\times})} \mc{O}_X^{\times}. \]
    The map to \(\mc{O}_X\) is given by \((a,b) \to \alpha(a) b\). 
\end{defn}

\begin{defn}
    A morphism of log schemes \((X, M_X) \to (Y, M_Y)\) consists of
    \begin{enumerate}
        \item A morphism \(\ul{f} \colon \ul{X} \to \ul{Y}\) of underlying schemes;
        \item A map \(f^{\flat} \colon f^* M_Y \to M_X\), where \(f^* M_Y\) is the log structure associated to \(f^{-1} M_Y\).  
    \end{enumerate}
\end{defn}

\begin{defn}
    A \textit{chart} for a log scheme \((X, \alpha \colon M_X \to \mc{O}_X)\) is a morphism from the constant monoid \(P_X \to M_X\) such that \(M_X \simeq P_X^a \simeq P_X \oplus_{\alpha_P^{-1}(\mc{O}_X^{\times})} \mc{O}_X^{\times}\), where \(\alpha_P \colon P_X \to M_X \to \mc{O}_X\).
\end{defn}

Charts are useful in specifying log structures. For example, if we have \(P_X \to \mc{O}_X\), then \(P_X^a\) gives a log structure on \(X\).

\begin{exm}
    The trivial log structure is given by the chart \(P = \ab\{0\}\).
\end{exm}

\begin{exm}
    Let \(X = \Spec k\) and give a log structure by \(\alpha_{\N} \colon \N \to k\) which sends \(0\) to \(1\) and \(1\) to \(0\). This is called the \textit{standard log point} and can also be formed by restricting the divisorial log structure on \(\A^1\) to the origin.
\end{exm}

\begin{exm}
    Consider the chart \(\N \to \mc{O}_{\A^1}\) which sends \(n \mapsto t^n\). Then this recovers the divisorial log structure on \(\A^1\) relative to the origin.
\end{exm}

\begin{exm}
    Consider the chart \(\N^2 \to \mc{O}_{\A^2}\) which sends \((a,b) \mapsto x^a y^b\). This recovers the divisorial log structure on \(\A^2\) relative to the toric boundary.   
\end{exm}

\begin{defn}
    A \textit{chart} for a log map \(f \colon (X, M_X) \to (Y, M_Y)\) is a triple \((P_X\to M_X, Q_Y \to M_Y, Q \to P)\), where \(P_X \to M_X\) and \(Q_Y \to M_Y\) are charts and the diagram
    \begin{equation*}
    \begin{tikzcd}
        Q_X \ar{r} \ar{d} & P_X \ar{d} \\
        f^* M_Y \ar{r} & M_X
    \end{tikzcd}
    \end{equation*}
    commutes.
\end{defn}

\begin{defn}
    A log scheme \((X, M_X)\) is \textit{finitely generated, fine, integral, saturated} if there exists an \'etale local chart \(P_U \to M_U\) for a cover \(U \to X\) such that \(P\) is finitely generated, fine, integral, or saturated, respectively.
\end{defn}

\subsection{Log smoothness}%
\label{sub:Log smoothness}

\begin{defn}
    Consider a commutative diagram
    \begin{equation*}
    \begin{tikzcd}
        T_0 \ar{r}{\phi} \ar{d}{j} & X \ar{d}{f} \\
        T_1 \ar{r}{\psi} & Y,
    \end{tikzcd}
    \end{equation*}
    where \(j\) is a strict closed immersion (here, this means that \(M_{T_0} = j^* M_{T_1}\)). A morphism \(f \colon X \to Y\) of fine log schemes is \textit{log smooth (resp. log \'etale)} if
    \begin{enumerate}
        \item The underlying map \(\ul{f} \colon \ul{X} \to \ul{Y}\) is locally of finite presentation;
        \item For any commutative diagram as above, \'etale locally on \(T_1\) there exists a (resp. there exists a unique) morphism \(g \colon T_1 \to X\) such that \(\phi = g \circ j\) and \(\psi = f \circ g\).
    \end{enumerate}
\end{defn}

We will now give a chart local criterion for log smoothness.

\begin{thm}[{\cite[Theorem~3.5]{Kato1989}}]
    Let \(f \colon X \to Y\) be a morphism of fine log schemes. Assume we have a chart \(Q \to M_Y\) where \(Q\) is a fine monoid. Then \(f\) is log smooth if and only if \'etale locally on \(X\) there exists a chart for \(f\) extending \(Q \to M_Y\) satisfying
    \begin{enumerate}
        \item The kernel and the torsion part of the cokernel of \(Q^{\mr{gp}} \to P^{\mr{gp}}\) are finite groups whose orders are invertible on \(X\);
        \item The induced morphism \(\ul{X} \to \ul{Y} \times_{\Spec \Z[Q]} \Spec \Z[P]\) is \'etale in the classical sense.
    \end{enumerate}
\end{thm}

\begin{exm}
    Let \(X\) be an affine toric variety given by \(\Spec \C[P]\) for a monoid \(P\). Note that \(P^a \simeq M_{X, X \setminus T}\) induces the divisorial log structure. Then consider the map \(f \colon X \to \mr{pt}\), where the point has a trivial log structure. We note that the morphism \(0 \to P^{\mr{gp}}\) has no kernel and the cokernel is torsion-free, and the map
    \[ \ul{X} \to \Spec \C \times_{\Spec \Z} \Spec \Z[P] = \Spec \C[P] \]
    is an isomorphism, so any toric variety is log smooth over a point with trivial log structure.
\end{exm}

From this example, we can morally think of log smoothness as being locally toric relative to the base.

\begin{exm}
    Recall the family \(xy = t\) with log structure relative to the central fiber over \(\A^1\) with the divisorial log structure relative to \(0\). The map of charts we will consider is in fact \(\Delta \colon \N \to \N^2\). Taking the groupification, the map becomes \(\Delta \colon \Z \to \Z^2\). The kernel is trivial and the cokernel is torsion-free, so we have verified the first condition. Then we consider the map
    \[ \ul{X} \to \Spec \C[t] \times_{\Spec \Z[\N]} \Spec \Z[\N^2] = \Spec \frac{\C[u,v,t]}{uv-t}, \]
    which is an isomorphism, so this family is log smooth.
\end{exm}

\subsection{Stack of log structures}%
\label{sub:Stack of log structures}

The stack of log structures and its relation to log smoothness are due to Olsson~\cite{Olsson2003}.

Let \(S\) be a fine log scheme. We will consider the category \(\ms{Log}_S\) whose objects are fine log schemes \((T, M_T) \to S\) and whose morphisms are strict morphisms \((T, M_T) \to (T', M_{T'})\) over \(S\).
\begin{enumerate}
    \item This category \(\ms{Log}_S\) is a category fibered in groupoids over \(\ms{Sch}_{/S}\) classifying all fine log schemes mapping to \(S\);
    \item \(\ms{Log}_S\) is an algebraic stack locally of finite type over \(S\);
    \item A morphism \(f \colon X \to S\) is log smooth if and only if \(X \to \ms{Log}_S\) is smooth. The same holds with smooth replaced by \'etale.
\end{enumerate}

\subsection{Tropicalization}%
\label{sub:Tropicalization}

Let \(X\) be a fine log scheme. Then the tropicalization of \(X\) is defined to be the colimit
\[\Sigma(X) \coloneqq \bigsqcup_{x \in X} \Hom(\bar{M}_{X,x}, \R_{\geq 0}) \biggm/ \sim, \]
where the equivalence relation is given by if a point \(x\) specializes to \(y\), then the induced map \(\sigma_x \to \sigma_y\) induces an equivalence from \(\sigma_x\) to its image. 

Now suppose there is a stratification of \(X\) by \(X_i\) such that \(\bar{M}_{|X_i}\) is constant. For any \(X_i\), all points \(x \in X_i\) satisfy \(x \in \ol{\ab\{\eta_{X_i}\}}\), so we only need to consider the generic points of the strata and then glue various faces. 

\begin{exm}
    Consider \(X\) with the trivial log structure. Note that for any point \(x \in X\), we have \(\bar{M}_{X,x} = 0\), so \(\sigma_X = \ab\{0\}\). Then \(\Sigma(X) = \ab\{0\}\).
\end{exm}

\begin{exm}
    Consider \(\A^1\) relative to the origin (with divisorial log structure). It is easy to see that \(\sigma_0 = \Hom(\N, \R_{\geq 0}) = \R_{\geq 0}\) and \(\sigma_{\eta} = \ab\{0\}\), so the tropicalization is \(\R_{\geq 0}\).
\end{exm}

\begin{exm}
    Consider \(\A^2\) with the toric log structure (divisorial log structure relative to \(xy = 0\)). Then \(\sigma_{(0,0)} = \R^2_{\geq 0}\), \(\sigma_{(x,0)} = \R_{\geq 0}\), \(\sigma_{(0,y)} = \R_{\geq 0}\), and \(\sigma_{(x,y)} = \ab\{0\}\), so we obtain the fan of \(\A^2\) as a toric variety treated as a cone complex.
\end{exm}

\begin{exm}
    Let \(X_{\Sigma}\) be a toric variety with the divisorial log structure relative to \(X_{\Sigma} \setminus T\). Then
    \[ \Sigma(X_{\Sigma}) = \Sigma, \]
    so a possible slogan is that tropicalizing a toric variety recovers the structure of ordinary toric geometry (minus the embedding of the fan in some real vector space).
\end{exm}

\begin{exm}
    Let \(X = \Spec k[x,y,t]/(xy-t^{\ell})\) with divisorial log structure over \(xy = 0\). Then there is a chart given by
    \[ Q = \ab<e_x, e_y, e_t \mid e_x + e_y = \ell e_t>. \]
    Note that \(M_X = Q^a = Q \oplus_{\alpha^{-1}(\mc{O}_X^{\times})} \mc{O}_X^{\times}\) and \(\bar{M}_X = Q\). Therefore, we have
    \begin{align*}
        \sigma_q &= \Hom(Q, \R_{\geq 0}) \\
        &= \ab\{ (\rho_x, \rho_y, \rho_t) \in \R_{\geq 0}^3 \mid \rho_x + \rho_y = \ell \rho_t \} \\
        &= \ab\{ (\rho_t, \rho_x) \in \R_{\geq 0}^2 \mid 0 \leq \rho_x \leq \ell \rho_t \}
    \end{align*}
    at the singular point \(q\), which recovers the cone which gives \(X\) as a toric variety. At the generic point \(\eta_x\) of \(x = 0\), we see that \(\bar{M}_{X, \eta_x} = \N e_t\), and therefore \(\sigma_{\eta_x} = \R_{\geq 0}\). Similarly, we see that \(\sigma_{\eta_y} = \R_{\geq 0}\). The remaining problem is gluing. 

    Note that \(\bar{M}_{X, q} \to \bar{M}_{X, \eta_x}\) is given by \(e_t \mapsto e_t\), \(e_x \mapsto \ell e_t\), and \(e_y \mapsto 0\), and therefore the map \(\sigma_{\eta_x} \to \sigma_q\) sends \(\rho_t\) to \((\rho_t, \rho_x)\), where \(\rho_x = \ell \rho_t\).

    The morphism \(X \to \A^1_t\) giving a family of curves with nodal central fiber has tropicalization given by projecting the cone onto \(\R_{\geq 0}\).
\end{exm}


\subsection{Log curves}%
\label{sub:Log curves}

For now, all monoids and log schemes will be fine (or even fine saturated). There are two reasons for this:
\begin{enumerate}
    \item If \(P\) is a fine monoid, its dual is fs and sharp;
    \item If \(X = (\ul{X}, M_X)\) is a fine log scheme, for any geometric point \(\bar{x} \to X\), there exists an \'etale neighborhood of \(\bar{x}\) which has a chart.
\end{enumerate}

We will now discuss various geometric properties of maps between monoids.
\begin{defn}
    A map \(\theta \colon P \to Q\) is \textit{integral} if the map \(\Spec\C[Q] \to \Spec \C[P]\) is flat.
\end{defn}

For example, any nonzero map \(\N \to \N\) and the diagonal \(\Delta \colon\N \to \N^2\) are integral. If a map is integral, then for any \(P \to P'\), the monoid \(P' \oplus_P Q\) is fine.

\begin{exm}
    Pushouts may not be fine, for example consider
    \[ \ab< a ,b,c \mid a+c=b+c> = \N^3 \oplus_{\N^2} \N. \]
\end{exm}

\begin{defn}
    A map \(f \colon X \to Y\) between fine log schemes is integral if at every geometric point of \(X\), the map
    \[ (f^*\bar{M}_Y)_{\bar{x}} \to \bar{M}_{X,\bar{x}} \]
    is integral.
\end{defn}

\begin{defn}
    A map between fs monoids \(P \to Q\) is \textit{saturated} if for all \(P \to P'\), the monoid \(P' \oplus_P Q\) is also fs.
\end{defn}

\begin{exm}
    The map \(\Delta \colon \N \to \N^2\) is saturated, but \(\ell \colon \N \to \N\) is not.
\end{exm}

\begin{defn}
    A \textit{log curve} over an fs log scheme \(S\) is a log smooth and integral map \(\pi \colon \mc{C} \to S\) between fs log schemes such that the underlying geometric fibers are reduced connected curves.
\end{defn}

\begin{prop}[{\cite{Kato1989}}]
    Let \(f \colon X \to Y\) be a log smooth and integral map between fs log schemes. Then \(\ul{f}\) is flat.
\end{prop}

\begin{thm}[{\cite{FKato2000}}]
    The underlying family of a log curve is a family of nodal curves.
\end{thm}

\begin{exm}
    Let \(P = \N \setminus \ab\{1\}\). Then \(\Spec (P \to k[P])\) is log smooth over \(k\), but \(\Spec k[P]\) is a cusp.
\end{exm}

\begin{defn}
    An \textit{\(n\)-pointed log curve} over an fs base \(S\) consists of the data \((\pi \colon \mc{C} \to S), \vec{p} = \ab\{p_i\}_{i=1}^n\) such that
    \begin{enumerate}
        \item The underlying family is a prestable curve;
        \item \(\pi\) is a log smooth map between fs log schemes;
        \item On the smooth locus \(U\) of \(\ul{\pi}\), we have
            \[ \bar{M}_{\mc{C}}|_U \simeq \pi^* \bar{M}_S \oplus \bigoplus_{i=1}^n p_i \ul{\N}. \]
    \end{enumerate}
\end{defn}

